\begin{summarybox}
	\textbf{\textcolor{CSummary}{一句话核心}}

	\textbf{从 $y=\sin x$ 到 $y=A\sin(\omega x+\varphi)$，调整平移量后，两种水平变换顺序得到同一图像。}

	\medskip
	\textbf{\textcolor{CSummary}{条件与速查}}
	\begin{itemize}[leftmargin=*, itemsep=0.5em, topsep=0.4em]
		\item $A\neq0$，$\omega>0$，$\varphi\in\mathbb R$；周期 $2\pi/\omega$，振幅 $|A|$。
		\item 先向左平移有符号距离 $\varphi$，再将横坐标乘以 $1/\omega$；或先横向伸缩，再向左平移 $\varphi/\omega$。两条路径最后都将纵坐标乘以 $A$。
		\item $\varphi>0$ 左移，$\varphi<0$ 右移，$\varphi=0$ 不移动；两段距离分别为 $|\varphi|$、$|\varphi|/\omega$，当且仅当 $\varphi=0$ 或 $\omega=1$ 时相等。
		\item $A<0$ 时，纵坐标乘以 $A$ 等价于先按 $|A|$ 倍伸缩再关于 $x$ 轴对称；$\omega>1$ 横向压缩，$\omega=1$ 不变，$0<\omega<1$ 横向拉伸。
	\end{itemize}

	\textbf{\textcolor{CSummary}{最后检查}}
	\textbf{比较同一点的最终横坐标：$(t-\varphi)/\omega=t/\omega-\varphi/\omega$。等价来自平移量的调整，并非两种操作一般可交换。}
\end{summarybox}
